\chapter*{Abstract}
\addcontentsline{toc}{chapter}{Abstract}

Colorectal cancer is among the most prevalent and deadly malignancies worldwide.
Early detection of somatic mutations in cell-free DNA (cfDNA) holds the promise
of non-invasive screening, yet distinguishing tumour-derived fragments from the
overwhelming background of healthy DNA remains a formidable challenge.

This thesis investigates a novelty detection approach based on string kernels
applied to short DNA reads. We represent each read in a high-dimensional feature
space induced by the \emph{mismatch string kernel}, which counts all $k$-mers
within a specified Hamming distance of the observed subsequences. An extended
six-letter epigenetic alphabet --- incorporating 5-methylcytosine (M) and
5-hydroxymethylcytosine (H) alongside the canonical bases --- enables the kernel
to capture methylation signatures that are known to be altered in cancer.

A One-Class Support Vector Machine (OC-SVM) is trained exclusively on reads from
healthy donors to learn a decision boundary that encloses the ``normal'' region
of the kernel-induced feature space. At inference time, each read from a test
patient receives an anomaly score; these scores are then aggregated at the
patient level following a Multiple Instance Learning (MIL) paradigm, treating
every patient as a \emph{bag} of read-level instances.

We evaluate the pipeline through a rigorous Leave-One-Out cross-validation
protocol over a cohort of healthy and colon-cancer patients, reporting
patient-level ROC-AUC as the primary metric. Experimental results demonstrate
that the mismatch kernel, combined with distance-weighted neighbourhood
expansion and Multiple Kernel Learning (MKL), yields a strong separation between
healthy and tumour patients across multiple random seeds.

The entire pipeline --- from memory-mapped FASTA parsing and parallel sparse
feature extraction to block-partitioned Gram matrix computation --- is
implemented in Python with an emphasis on computational efficiency and
reproducibility.
